% TOP_PROBLEM: {"id":30007729,"problem_number":"LOCAL-30007729","title":"Realized returns to medical AI","category_id":21,"category":{"id":21,"name":"economics","display_name":"Economics","description":"Open economic theory statements, published research agendas, and proposed scoped specifications; question provenance and historical priority are qualified per record.","slug":"economics","order_index":21,"created_at":"2026-10-08T00:00:00Z"},"status":"provenance_pending","external_url":null,"legacy_ids":[],"published":false,"statement":"Estimate the realized two-year health, resource-cost, and clinician-labor effects of one version-controlled medical AI system in randomized clinics.","economics":{"original_id":218,"field":"Health economics","kind":"empirical","formalization_basis":"proposed_specification","source_status":"not_verified","priority_status":"not_established","priority_note":"No readable primary passage explicitly stating this precise open question was verified. A supporting paper, abstract, or topic citation is insufficient to establish origin or unresolved status.","earliest_verified_reference":null,"current_reference":{"authors":["Nikhil Sahni","George Stein","Rodney Zemmel","David M. Cutler"],"title":"The Potential Impact of Artificial Intelligence on Healthcare Spending","year":2023,"url":"https://www.nber.org/papers/w30857","doi":"10.3386/w30857","locator":"Abstract","evidence_summary":"The source estimates potential spending savings; projections do not establish the specified causal implementation question as an explicitly stated open problem."},"formalization_note":"Proposed scoped research specification. The original catalogue prompt was a synthesis, and no canonical conjecture or verified original formulation is claimed. The supporting source, when retained, establishes context or existing results rather than this exact open question."},"statusQualification":"Provenance pending: an explicit published statement of this exact open question was not verified. This is a catalogue-authored candidate specification, not a certified open conjecture. Proposed scoped research specification. The original catalogue prompt was a synthesis, and no canonical conjecture or verified original formulation is claimed. The supporting source, when retained, establishes context or existing results rather than this exact open question.","statusReviewedAt":"2026-10-08"}
% FIRST_POSTED: 2026-10-08
% ENTRY_KIND: statement_only
% !TeX program = lualatex
\documentclass[11pt]{article}
\usepackage{fontspec}
\usepackage[a4paper,margin=25mm]{geometry}
\usepackage{amsmath,amssymb,mathtools}
\usepackage{xurl}
\usepackage[hidelinks]{hyperref}
\newcommand{\sref}[2]{\hyperlink{source:#1:#2}{[S#2]}}
\setlength{\emergencystretch}{3em}
\begin{document}
% problemId:30007729
\section*{Realized returns to medical AI}\label{30007729}
\subsection{Definitions and mathematical statement}
Consider adult primary-care visits in a prespecified national provider network over two years. Randomize clinics to access \(a=1\) to one version-controlled artificial-intelligence system for documentation and decision support, or \(a=0\) to existing practice. Let \(Q_i(a)\) be quality-adjusted health, \(C_i(a)\) total real care costs, and \(L_i(a)\) clinician labor hours associated with patient \(i\). For health valuation \(\lambda>0\), define
\[\tau_W=E[\lambda\{Q_i(1)-Q_i(0)\}-\{C_i(1)-C_i(0)\}]-k,\qquad\tau_L=E[L_i(1)-L_i(0)],\]
where \(k\) is implementation, maintenance, supervision, and training resource cost per eligible patient. Costs include follow-up testing, downstream hospital use, and human review of errors. Determine realized welfare and labor effects under randomized clinic assignment, not savings implied by automatable task shares. Track uptake, diagnostic mistakes, altered visit volume, and effects by patient complexity. Preserve baseline patient inclusion and specify clinic spillovers and model-version changes. Longer-term organizational responses require a separate extension or explicit extrapolation assumptions. The cited spending paper projects potential gains; an explicit open statement of this causal implementation target was not verified.

\par\textbf{Formulation and scope.} Proposed scoped research specification. The original catalogue prompt was a synthesis, and no canonical conjecture or verified original formulation is claimed. The supporting source, when retained, establishes context or existing results rather than this exact open question.

\par\textbf{Open-question provenance.} An explicit published open-question statement was not verified. Sources below are supporting literature and must not be treated as a first listing.

\par\textbf{Historical priority.} No readable primary passage explicitly stating this precise open question was verified. A supporting paper, abstract, or topic citation is insufficient to establish origin or unresolved status.

\subsection{Short English statement}
Estimate the realized two-year health, resource-cost, and clinician-labor effects of one version-controlled medical AI system in randomized clinics.

\subsection{Sources}
\begin{itemize}\raggedright
\item[\textnormal{[S1]}]\hypertarget{source:30007729:1}{} Nikhil Sahni, George Stein, Rodney Zemmel, David M. Cutler. 2023. The Potential Impact of Artificial Intelligence on Healthcare Spending. Abstract.
\url{https://www.nber.org/papers/w30857}.
DOI: 10.3386/w30857.
{\small\textit{Supporting or current literature; see evidence qualification. The source estimates potential spending savings; projections do not establish the specified causal implementation question as an explicitly stated open problem.}}
\end{itemize}
\end{document}
